\section{Conclusion}

\begin{comment}
This project proposal is an outline of a multi-level, multi-protective architecture to reduce control-plane timing attacks in Software-Defined Vehicular Networks. The proposed solution is a new concept of enhancing the security, reliability and performance in dynamic vehicles environments, with a synergistic integration of blockchain technology, deep reinforcement learning and large language models. The framework particularly focuses on some of the major gaps observed in the modern studies. First, it has built-in defense, instead of single detection techniques, which combines tamper-proof timing-event recording with blockchain, adaptive anomaly detection and response provided by deep reinforcement learning, and smart threat analysis offered by language models on a large scale. This compound approach has the benefit of providing strong defense against multi-vector timing attacks that are highly complex and difficult to defend against under single-layered security defense.

Second, the framework puts primary emphasis on proactive defense, as opposed to reactive detection. When using such a strategy as deep reinforcement learning, the system can predict any possible attacks and adjust timing policies dynamically, which will avoid causing significant harm before its occurrence. This preemptive capability is of the utmost importance with regard to time critical safety applications to vehicular networks where delays of a single second can give rise to catastrophic effects.

Thirdly, the solution advanced is designed in a way that supporting the low-latency needs that cannot be afforded in the automobile sector is maintained and at the same time security measures are instituted. The framework attempts to maintain the network performance and time within a prescribed error margin and, at the same time, offers high security assurance through intelligent optimisation of defence activation and selective deployment of security functions.

The study is relevant to the development of the intelligent transportation systems security in the sense that it demonstrates how emerging technologies can be optimally integrated to address sophisticated security threats. The blockchain component guarantees the integrity of timing by immutable logging and verification based on consensus; the deep reinforcement learning component helps to adaptively learn and optimise timing policies; and the large language model component provides contextual knowledge and intelligent decision support in threat analysis.

It is expected that the improvements in the rate of attack detection with few false positives will be significant, the network performance will remain unaffected by both the normal and adversarial conditions, the timing accuracy and the quality of synchronization on the control plane would increase, and the practical implications of applying blockchain, deep reinforcement learning, and large language models to network security issues would be provided.

The proposed simulation-based analysis that will be used Network Simulator 3 will provide strict verification of the proposed ideas in a variety of attacks and network parameters. The effectiveness and benefits of the combined method will be proved by comparing performance analysis before and after implementation with solutions that have been studied successfully. The findings will be shared in peer-reviewed journal articles and conference papers, thus contributing to the body of knowledge and state-of-the-art in vehicles network security.

The project represents a major step towards securing the Software-Defined Vehicular Networks to control-plane timing attacks, and future usage can be seen to other critical infrastructure networks that require very precise timing and real-time communication. The fruitful completion of this study will provide theoretical underpinnings and practical guidelines on the future activities in intelligent, adaptive network security systems.
\end{comment}

This report has presented DMAP-SDVN, a dual-mode adaptive protection
framework for control-plane timing attacks in Software-Defined Vehicular
Networks. The framework addresses four attack variants — delay injection,
race-condition exploitation, timing side-channel leakage, and
synchronisation exploits — through a formally defined combination of
mobility-aware attack signatures, a DRL-based detection agent, a
Hyperledger Fabric blockchain layer with IPFS-anchored timing proofs,
a post-quantum cryptographic layer using ML-DSA and ML-KEM, and an
LLM-based contextual intelligence module. The dual-mode architecture
switches between a lightweight rule-based mode and a full AI/blockchain/PQC
mode based on real-time threat scores and resource availability,
addressing the practical deployment constraints of resource-limited RSUs.
Simulation-based evaluation using NS3 and SUMO against three external
baselines and five ablation variants is ongoing, with results to be
reported in the final submission.